% ============================================================================
% PLATZHALTER -- Bild 5: Entrollter Integrand, Bromwich vs. Talbot
% ============================================================================
% REGIEANWEISUNG FUER DIE GRAFIK-TOOLCHAIN (bevorzugt matplotlib/pgfplots,
% als PDF eingebunden oder direkt pgfplots; Schrift Times 10pt wie Text):
%
% DATENGRUNDLAGE (echte Berechnung, F(s)=1/s, t=1):
% - BROMWICH: s = gamma + i*omega mit gamma=1; Integrand
%   Re[ e^{st} F(s) ] = Re[ e^{t(gamma+i*omega)} / (gamma+i*omega) ]
%   als Funktion des Kurvenparameters omega in [-40, 40].
%   -> oszillierende Kurve mit nur ~1/|omega| abfallender Amplitude.
% - TALBOT: Kontur s(theta) = sigma + lambda(theta*cot(theta) + i*beta*theta)
%   mit z.B. sigma=0, lambda=2, beta=1, theta in (-pi, pi);
%   (beta ist der Formparameter, im Paper-Text s_beta; in der
%   Originalliteratur oft nu genannt -- hier NICHT nu beschriften!)
%   Integrand Re[ e^{t s(theta)} F(s(theta)) s'(theta) / (2*pi*i) ]
%   als Funktion von theta.
%   -> glatte, glockenfoermige Kurve, an den Raendern theta -> +-pi
%      praktisch null (superexponentieller Abfall).
%
% DARSTELLUNG:
% - ZWEI Panels nebeneinander (gleiche Hoehe), oder ein Panel mit zwei
%   Kurven, falls die Skalen vertraeglich sind; bevorzugt ZWEI PANELS:
%   * links:  "Bromwich-Gerade", Kurve in ROT (plotred), x-Achse
%     "$\omega$", y-Achse "Integrand".
%   * rechts: "Talbot-Kontur", Kurve in BLAU (plotblue), x-Achse
%     "$\theta$", y-Achse identisch skaliert beschriftet.
% - Beide y-Achsen mit derselben Einheit; falls die Amplituden stark
%   differieren, y-Limits pro Panel waehlen, aber im Caption-Text bleibt
%   die Aussage qualitativ.
% - Dezentes Gitter (grau, duenn), keine Titelzeile im Bild (Caption
%   uebernimmt das), Legenden nur falls beide Kurven in einem Panel.
% - Auf der Talbot-Kurve zusaetzlich N=12 Stuetzstellen der Trapezregel
%   als Punkte markieren (blaue Punkte), um zu zeigen, wie wenige
%   Auswertungen genuegen.
%
% GROESSE: Gesamtbreite ca. 12cm, Hoehe ca. 4.5cm.
% ============================================================================
\scalebox{0.8}{
\begin{tikzpicture}

% ==============================
% 1. FARBEN DEFINIEREN
% ==============================
\definecolor{plotblue}{RGB}{0, 0, 200}
\definecolor{plotred}{RGB}{220, 30, 30}
\definecolor{plotgray}{RGB}{140, 140, 140}
\definecolor{plotbg}{RGB}{244, 244, 252}

% ==============================
% 2. PGFPLOTS SETUP
% ==============================
\pgfplotsset{table/search path={.,Latex/images,images,papers/laplace/images}}

% ==============================
% 3. BROMWICH PANEL (Links)
% ==============================
\begin{axis}[
    name=ax1,
    width=7cm,
    height=5.5cm,
    grid=major,
    grid style={plotgray!30, thin},
    axis lines=left,
    axis line style={thick},
    tick label style={font=\small},
    label style={font=\normalsize},
    enlarge x limits=false,
    xlabel={$\omega$},
    ylabel={Integrand},
    xmin=-40, xmax=40,
    ymin=-0.8, ymax=3.0,
    title={\textbf{Bromwich-Gerade}},
    title style={font=\normalsize, yshift=-1.5ex}
]
\addplot[
    thick, 
    plotred,
    ] table[x=omega, y=integrand, col sep=comma] {papers/laplace/images/data_integrand_bromwich.csv};
\end{axis}

% ==============================
% 3. TALBOT PANEL (Rechts)
% ==============================
\begin{axis}[
    name=ax2,
    at={(ax1.outer east)},
    anchor=outer west,
    xshift=1.5cm,
    width=7cm,
    height=5.5cm,
    grid=major,
    grid style={plotgray!30, thin},
    axis lines=left,
    axis line style={thick},
    tick label style={font=\small},
    label style={font=\normalsize},
    enlarge x limits=false,
    xlabel={$\theta$},
    xmin=-3.1415, xmax=3.1415,
    ymin=-0.3, ymax=1.3,
    xtick={-3.1415, -1.5708, 0, 1.5708, 3.1415},
    xticklabels={$-\pi$, $-\pi/2$, $0$, $\pi/2$, $\pi$},
    title={\textbf{Talbot-Kontur}},
    title style={font=\normalsize, yshift=-1.5ex}
]
\addplot[
    thick, 
    plotblue,
    ] table[x=theta, y=integrand, col sep=comma] {papers/laplace/images/data_integrand_talbot.csv};

\addplot[
    only marks,
    mark=*,
    mark size=1.5pt,
    plotblue,
    ] table[x=theta, y=integrand, col sep=comma] {papers/laplace/images/data_integrand_talbot_pts.csv};
\end{axis}

\end{tikzpicture}
}
